% Corrected French for printed p.186; departures are append-only below.
\clearpage
\SourcePage{186}

\LayerTarget{tohoku:proof:proposition-4.1.1:p1}
\LayerTarget{tohoku:diagram:4.1.1:finite-type-stalk-hom}
\[
\begin{array}{ccccc}
0&\longrightarrow&\mathbf{Hom}_O(A,B)(x)&\longrightarrow&
  \mathbf{Hom}_O(O^n,B)(x)\\
&&\big\downarrow&&\big\downarrow\\[0.45em]
0&\longrightarrow&\mathrm{Hom}_{O(x)}(A(x),B(x))&\longrightarrow&
  \mathrm{Hom}_{O(x)}(O^n(x),B(x)).
\end{array}
\]

\LayerTarget{tohoku:eq:4.1.4}
Or on a $\mathbf{Hom}_O(O^n,B)=\mathbf{Hom}_O(O,B)^n=B^n$, car on a un
isomorphisme canonique
\[
\tag{4.1.4}
\mathbf{Hom}_O(O,B)=B.
\]

\LayerTarget{tohoku:proof:proposition-4.1.1:p2}
Les groupes dans la dernière colonne du diagramme s'identifient donc
respectivement à $B(x)^n$ et
$\mathrm{Hom}_{O(x)}(O(x),B(x))^n=B(x)^n$, donc le dernier homomorphisme
vertical est un isomorphisme, d'où on conclut aussitôt que le premier
homomorphisme est un monomorphisme. Supposons maintenant que $A$ soit
pseudo-cohérent en $x$, c'est-à-dire (en se restreignant au besoin à un
voisinage convenable de $x$) qu'on a une suite exacte
$O^m\longrightarrow O^n\longrightarrow A\longrightarrow0$. Se servant
toujours de l'exactitude à gauche des foncteurs $\mathrm{Hom}$, on en conclut
encore un homomorphisme de suites exactes :

\LayerTarget{tohoku:diagram:4.1.1:pseudo-coherent-stalk-hom}
\[
\begin{array}{ccccccc}
0&\longrightarrow&\mathbf{Hom}_O(A,B)(x)&\longrightarrow&
  \mathbf{Hom}_O(O^n,B)(x)&\longrightarrow&\mathbf{Hom}_O(O^m,B)(x)\\
&&\big\downarrow&&\big\downarrow&&\big\downarrow\\[0.45em]
0&\longrightarrow&\mathrm{Hom}_{O(x)}(A(x),B(x))&\longrightarrow&
  \mathrm{Hom}_{O(x)}(O^n(x),B(x))&\longrightarrow&
  \mathrm{Hom}_{O(x)}(O^m(x),B(x)).
\end{array}
\]

\LayerTarget{tohoku:proof:proposition-4.1.1:p3}
Nous avons déjà remarqué que les deux derniers homomorphismes verticaux sont
bijectifs. Il en résulte aussitôt que le premier homomorphisme vertical est un
isomorphisme.

\LayerTarget{tohoku:question:injectivity-of-stalk}
Supposons maintenant que $B$ soit un objet injectif de la catégorie abélienne
$C^O$, peut-on en conclure que $B(x)$ est un $O(x)$-module injectif ? En vertu
du lemme 1 de 1.10., il suffit de vérifier que pour tout idéal à gauche $a$ de
$O(x)$ et tout $O(x)$-homomorphisme de $a$ dans $B(x)$, ce dernier peut se
prolonger en un $O(x)$-homomorphisme de $O(x)$ dans $B(x)$. Comme la
restriction d'un $O$-Module injectif à un ouvert est encore injective (prop.
3.1.3), il suffit pour ceci que l'on puisse trouver un voisinage ouvert $U$ de
$x$, un sous-Module $A$ de $O|U$, enfin un homomorphisme $u$ de $A$ dans
$B|U$, de telle façon que $A(x)=a$ et que l'homomorphisme de $A(x)$ dans
$B(x)$ défini par $u$ soit l'homomorphisme donné $a\longrightarrow B(x)$.
D'ailleurs, si $A$ est pseudo-cohérent en $x$, l'existence de $u$ résulte
automatiquement de prop. 4.1.1. Prouvons alors :

\LayerTarget{tohoku:lemma:4.1.1}
\noindent\textsc{Lemme 1.} \textit{Supposons que $O$ soit un faisceau cohérent
d'anneaux [15, Chap. I, par. 2], soit $M$ un $O$-Module cohérent (``faisceau
cohérent de modules'' dans la terminologie de [15]), et $n$ un sous-module de
type fini de $M(x)$. Alors il existe, sur un voisinage ouvert convenable $U$ de
$x$, un sous-faisceau cohérent $N$ de $M$ tel que $N(x)=n$.}

\LayerTarget{tohoku:proof:lemma-4.1.1}
Soit $(n_i)$ un système fini de générateurs de $n$ ($1\leq i\leq k$). Pour
tout $i$, on a $n_i\in M(x)$, donc $n_i=f_i(x)$, où $f_i$ est une section de
$M$ définie dans un voisinage ouvert de $x$. On peut supposer tous les $f_i$
définis dans un même voisinage ouvert $U$, ils définissent alors au-dessus de
$U$ un homomorphisme $O^k\longrightarrow M$. On prendra pour $N$ l'image de
cet homomorphisme, qui est cohérent comme sous-module de type fini du faisceau
cohérent $M$.

\TohokuContinuation{187}
